% !TEX program = lualatex
\documentclass[12pt,a4paper]{article}
\usepackage{luatexja}
\usepackage{amsmath,amssymb,amsthm}
\usepackage{bm}
\usepackage{hyperref}
\usepackage{geometry}
\geometry{margin=1in}
\usepackage{enumitem}

%-------------------------------------------------
%  マクロ
%-------------------------------------------------
\newcommand{\dfix}{D_{\mathrm{fix}0}} % 0次元不動点（空）
\newcommand{\Rep}{\mathcal{R}}
\newcommand{\Prf}{\mathcal{P}}
\newcommand{\Tset}{\mathcal{T}}

%-------------------------------------------------
%  定理環境
%-------------------------------------------------
\newtheorem{definition}{定義}[section]
\newtheorem{theorem}{定理}[section]
\newtheorem{remark}{注}[section]

\begin{document}

\title{正当性原理（Justification Principle）メモ}
\author{}
\date{}
\maketitle

\section{正当性原理（Justification Principle）}
本メモでは、\emph{正当性原理}を次の機序として捉える：
命題（静的情報）に対して\textbf{時間軸を一時的に付加}し、その後\textbf{時間軸を除去}することで観測者を固定し、結果の再現性を最大化する。

この原理は、
「不動点に不動点を入力しても常に同じ不動点が返ってくる」という強い閉包性（固定点の安定性）を持ち、
「正しさの判定そのものが、さらに正しさとして判定可能である」という性質を本質とする。

\section{証明とは何か}
\begin{remark}
証明とは、命題に対して正当性原理を適用する行為である。
すなわち（亦是亦非）や（非是非非）に対応する要素（＝時間軸に載った揺らぎ／不整合性）をいったん除去し、
正当なもののみを復元して「成立するか」を問う。
\end{remark}

\section{証明という行為の動機}
2 次元（静的）情報 $2d\setminus T$ に対する数学の証明行為は、
時間軸を付加しつつ最終的に除去することで観測者を固定し、
結果の再現性を最大化することを根本的な動機として持つ、と整理できる。

\section{2^n d に拡張した「証明」の模型}
以下では、$2d$ の議論を $(2^n)d$ へ拡張した形のメモをまとめる。

\begin{definition}[Proof in $(2^{n})d$]
$\Sigma$ を記号列（ビット列等）の集合、$\Tset$ を時刻トークンの集合とする。
\[
  \Rep_{n} := \{(s_{1},\dots ,s_{2^{n}})\mid s_{i}\in\Sigma\},
  \qquad
  \Prf_{n} := \{((s_{1},t_{1}),\dots ,(s_{2^{n}},t_{2^{n}}))\mid s_i\in\Sigma,\ t_{i}\in\Tset\}.
\]
埋め込み（時間軸付加）を
\[
  \iota_{n}\colon\Rep_{n}\hookrightarrow\Prf_{n},
  \qquad
  \iota_{n}(s_{1},\dots ,s_{2^{n}})
    = ((s_{1},t_{1}),\dots ,(s_{2^{n}},t_{2^{n}})),
\]
評価（射影、時間軸除去）を
\[
  \phi_{n}\colon\Prf_{n}\twoheadrightarrow\Rep_{n},
  \qquad
  \phi_{n}(((s_{1},t_{1}),\dots ,(s_{2^{n}},t_{2^{n}})))
    = (s_{1},\dots ,s_{2^{n}})
\]
で定める。
距離を
\[
  \|X\|:=K(X)+\ell(X),\qquad
  \ell(\Rep_{n})=0,\ \ell(\Prf_{n})=1
\]
とする。
\end{definition}

\begin{theorem}[Convergence in $(2^{n})d$]
任意の $p\in\Rep_{n}$ に対し列
\[
  v_{0}=p,\qquad v_{k+1}:=\phi_{n}\bigl(\iota_{n}(v_{k})\bigr)
\]
を定める。
さらに各ステップで距離が 1 だけ減少すると仮定する：
\[
  \|v_{k+1}\|=\|v_{k}\|-1\quad(k\ge0).
\]
このとき
\[
  \lim_{k\to\infty}v_{k}=\dfix,
\]
ただし $\dfix$ は $\|\dfix\|=0$ を満たす一意な点（$0$ 次元不動点）である。
\end{theorem}

\section{正しさの機序と LLM の幻覚（メモ）}
LLM の幻覚（hallucination）は「コスト最小化（局所最適化）」の結果として現れる現象であり、
これを抑止するにはコスト最小化とは別の\textbf{適切な判定関数}を確立する必要がある、と一般化できる。

\end{document}
